% Introduction for EMNLP - Temporal OOD focus

Knowledge graph (KG) embeddings underpin many downstream applications.
Question answering systems retrieve facts from Freebase and Wikidata~\citep{bordes2015large}, dialogue agents ground responses in structured knowledge~\citep{moon2019opendialkg}, and fact verification systems check claims against knowledge bases~\citep{thorne2018fever}.
The reliability of these applications depends on knowing \emph{when KG predictions are trustworthy}---yet uncertainty quantification for KG embeddings remains underexplored.

\paragraph{The Temporal Shift Problem.}
Unlike static benchmarks, real-world KGs evolve continuously, creating two distinct OOD scenarios that require different detection strategies:
\begin{enumerate}
    \item \textbf{Emerging entities}: New actors, products, or concepts appear that were absent during training (e.g., a newly elected official, a recent scientific discovery).
    \item \textbf{Novel relational contexts}: Familiar entities appear in previously unobserved relationships (e.g., a well-known scientist winning a new award type, a company entering a new market).
\end{enumerate}
NLP systems must detect \emph{both} failure modes and abstain rather than hallucinate.
This motivates \textbf{uncertainty-aware KG reasoning} that can detect temporal distribution shifts---but as we show, existing methods conflate these two scenarios and fail on both.

\paragraph{Scope.}
This paper addresses uncertainty quantification for KG embeddings under temporal distribution shift.
We evaluate on standard KG benchmarks (FB15k-237, WN18RR, YAGO3-10) and a temporally-stamped dataset (ICEWS14), with a selective prediction task to demonstrate practical utility.
Integration with end-to-end NLP systems (e.g., KGQA, dialogue) is future work; our contribution is the uncertainty framework for KG embeddings.

\paragraph{Why Current Methods Fail.}
Probabilistic KG embeddings~\citep{Chen2021PERM,chen2021probabilistic} learn a variance $\sigma^2_e$ for each entity, capturing uncertainty about entity representations.
The intuition is sound: rare entities should have high variance.

However, this approach has a fundamental limitation: learned variance is \emph{relation-agnostic}.
Consider an entity $e$ appearing frequently with \texttt{born\_in} but never with \texttt{works\_at}.
A single $\sigma^2_e$ cannot distinguish high-confidence $(e, \texttt{born\_in}, ?)$ from uncertain $(e, \texttt{works\_at}, ?)$.
Under temporal shift, this conflation causes systematic failures.

\paragraph{Our Solution: Semantic-Structural Decomposition.}
We show that uncertainty decomposes into two orthogonal components, each targeting one of the temporal OOD scenarios:
\begin{itemize}
    \item \textbf{Semantic uncertainty}: How well is this entity learned? (from embedding variance)---detects \emph{emerging entities}
    \item \textbf{Structural uncertainty}: Have we seen this entity-relation pair? (from coverage)---detects \emph{novel relational contexts}
\end{itemize}

These signals are complementary: semantic uncertainty excels on emerging entities (0.83 AUROC) over structural (0.78), while structural uncertainty achieves near-perfect detection on novel contexts (1.00) where semantic fails (0.42).
Neither alone suffices; combining both with adaptive weights achieves robust detection across all temporal shift scenarios.

\paragraph{Contributions.}
Our primary contribution is \emph{analytical}: we identify and formalize a fundamental limitation in existing uncertainty methods, then provide a principled decomposition framework. The algorithmic instantiation (CAGP) is deliberately simple---our claim is not algorithmic novelty but rather that \textbf{understanding the structure of uncertainty} enables effective solutions.

\begin{enumerate}
    \item \textbf{Analysis}: We show that all entity-level variance methods (GP-KGE, BEUrRE, UKGE) are relation-agnostic, explaining their systematic failure on temporal OOD (\S\ref{sec:method}).

    \item \textbf{Framework}: We formalize the semantic-structural decomposition and empirically demonstrate complementarity---each signal captures failure modes the other misses (\S\ref{sec:method}).

    \item \textbf{Empirical validation}: Score-based methods achieving 0.99 AUROC on standard benchmarks drop to 0.52--0.55 on temporal OOD; our framework achieves 0.87--0.97 across FB15k-237, WN18RR, YAGO3-10, and ICEWS14 with real timestamps (\S\ref{sec:experiments}).

    \item \textbf{Practical impact}: 48\% error reduction on selective link prediction at 85\% coverage (\S\ref{sec:experiments}).
\end{enumerate}

Our code and evaluation protocols are publicly available.\footnote{URL anonymized for review.}

Figure~\ref{fig:main_results} summarizes our main results: decomposition-based methods (CAGP, RelCondVar) achieve 0.87--0.97 AUROC across all datasets, substantially outperforming both single-signal baselines (GP-only, Coverage-only) and score-based methods (UKGE, Energy).
The gap is largest on temporal OOD, where score-based methods drop to near-random performance.

\begin{figure}[t]
\centering
\includegraphics[width=0.85\linewidth]{figures/fig1_main_results.pdf}
\vspace{-1em}
\caption{OOD detection AUROC across methods and datasets. Our decomposition-based approaches (CAGP variants) consistently outperform single-signal baselines. Error bars show $\pm$1 standard deviation across 3 seeds.}
\label{fig:main_results}
\end{figure}
